\documentclass[10pt,a4paper]{amsart}
\usepackage[margin=19mm]{geometry}
\usepackage{amsmath,amssymb,mathtools}
\usepackage[scr=rsfs,cal=euler]{mathalfa}
\usepackage{xfrac}
\usepackage[unicode,hidelinks,bookmarks=false]{hyperref}
\newcommand{\ie}{i.e.\ }
\newcommand{\cf}{cf.\ }
\newcommand{\resp}{respectively\ }
\newcommand{\Subsection}{\subsection}
\providecommand{\nobreakdash}{\nobreak\hbox{-}}
\setlength{\emergencystretch}{2em}
\multlinegap0pt
\usepackage{amssymb}
\usepackage{mathtools}
\usepackage{xcolor}
\usepackage{cancel}
\usepackage[shortlabels]{enumitem}
\newtheorem{corollary}{Corollary}
\newcommand{\C}{\mathbf{C}}
\newcommand{\I}{\mathbf{I}}
\newcommand{\bmu}{\boldsymbol{\mu}}
\newcommand{\bv}{\mathbf{v}}
\title{When Stronger Triggers Backfire:\\ A High-Dimensional Theory of Backdoor Attacks}
\author{Donald Flynn, Hadas Yaron Goldhirsh, Jonathan P. Keating, Inbar Seroussi}
\date{Source: arXiv:2605.22481v1 (CC BY 4.0)}
\begin{document}
\maketitle

\noindent\emph{Source and licence.} When Stronger Triggers Backfire:\\ A High-Dimensional Theory of Backdoor Attacks. Version arXiv:2605.22481v1 (CC BY 4.0), distributed by arXiv under the Creative Commons Attribution 4.0 International licence, http://creativecommons.org/licenses/by/4.0/, as recorded in the arXiv licence metadata for this exact version and shown on the abstract page https://arxiv.org/abs/2605.22481v1. Full licence notice: \texttt{licenses/arxiv-2605.22481-CC-BY-4.0.txt}.\par\smallskip

\noindent\textit{Editorial scope.} Counted unit: the article's corollary cor:lin\_peak\_eigen (source0.tex lines 530--555). The article's complete original proof of it is delivered with it (source0.tex lines 1792--2025, one \texttt{proof} environment of 234 lines, which the article places away from the statement). No mathematics is written, repaired or re-derived: the statement and the proof are the article's own lines, in the article's own notation, and no retained line is substituted or re-flowed.  Retained uncounted as the source-local context the proof needs: lines 1091--1101; lines 1103--1113; lines 1115--1123.  Nothing was omitted from any retained span, so the package carries no omission record.  No retained line contains a citation, so the wrapper carries no bibliography and no bibliography entry.  No retained line contains a figure environment, an image-inclusion command or drawing code, so no image is delivered and the package declares no asset dependency.  Editorial formatting only, in addition: an AMS \texttt{amsart} preamble that restates the article's own declarations and macros used by the retained lines, namely the environments \texttt{corollary} on separate counters, together with the standard packages those lines need.  The retained environments print in this document as Corollary 1 in order of appearance, whereas the article prints the same items with the numbering of its own full document.  The complete native source of the article as distributed in the arXiv e-print for this exact version is bundled unchanged as \texttt{sources/201140/source0.tex}.
\begin{equation}
\label{eq:lin_hmu_exact}
h_\mu(\alpha)
=
\frac{
\tau\Bigl((1-2\phi)g_{\mu\mu}
+\phi\alpha g_{\mu v}
+\tau\phi(1-\phi)\alpha^2\Delta_G\Bigr)
}
{D_{\mathrm{proj}}(\alpha)}
\end{equation}

\begin{equation}
\label{eq:lin_hv_exact}
h_v(\alpha)
=
\frac{
\tau\Bigl((1-2\phi)g_{\mu v}
+\phi\alpha g_{vv}
+2\tau\phi(1-\phi)\alpha\Delta_G\Bigr)
}
{D_{\mathrm{proj}}(\alpha)},
\end{equation}

\begin{equation}
\label{eq:lin_Dh}
D_{\mathrm{proj}}(\alpha)
=
1+\tau g_{\mu\mu}
-2\tau\phi\alpha g_{\mu v}
+\tau\phi\alpha^2 g_{vv}
+\tau^2\phi(1-\phi)\alpha^2\Delta_G .
\end{equation}

\begin{corollary}[Eigenvector specialisation]
\label{cor:lin_peak_eigen}
Assume  $\bmu^\top\bv=0$, $\C\bmu=s_\mu^2\bmu$, $\C\bv=s_v^2\bv$. Define $A_\mu:=\lambda+\tau s_\mu^2+\tau\|\bmu\|^2$, $B_v:=\lambda+\tau s_v^2$, $P_\mu:=\lambda+\tau s_\mu^2+\tau(1-\phi)\|\bmu\|^2$, $Q_\mu:=\lambda+\tau s_\mu^2+2\tau(1-\phi)\|\bmu\|^2$, and $D_{\mathrm{eig}}(\alpha):=A_\mu B_v+\tau\phi P_\mu\alpha^2$. Then
\begin{equation}
\label{eq:lin_ab_eigen_compact}
h_\mu(\alpha)
=
\|\bmu\|^2
\frac{\tau(1-2\phi)B_v+\tau^2\phi(1-\phi)\alpha^2}
{D_{\mathrm{eig}}(\alpha)},
\qquad
h_v(\alpha)
=
\frac{\tau\phi Q_\mu\alpha}
{D_{\mathrm{eig}}(\alpha)}.
\end{equation}
Consequently, \(h_\mu(\alpha)\) is strictly increasing for \(\alpha>0\),
while \(h_v(\alpha)>0\) for every \(\alpha>0\) and has a unique finite
maximizer at
\begin{equation}
\label{eq:lin_alpha_star_eigen_compact}
(\alpha_\star^{\mathrm{eig}})^2
=
(A_\mu B_v)/(\tau\phi P_\mu).
\end{equation}
\end{corollary}

\begin{proof}[Proof of Corollary~\ref{cor:lin_peak_eigen}]
Assume
\[
\bmu^\top\bv=0,
\qquad
\C\bmu=s_\mu^2\bmu,
\qquad
\C\bv=s_v^2\bv.
\]
Set
\[
L_\mu:=\lambda+\tau s_\mu^2,
\qquad
r:=\|\bmu\|^2.
\]
Then
\[
A_\mu=L_\mu+\tau r,
\qquad
B_v=\lambda+\tau s_v^2,
\]
and
\[
P_\mu=L_\mu+\tau(1-\phi)r,
\qquad
Q_\mu=L_\mu+2\tau(1-\phi)r.
\]
Since
\[
\mathbf{R}(\lambda,\tau)=(\lambda\I+\tau\C)^{-1},
\]
the eigenvector assumptions give
\[
\mathbf{R}(\lambda,\tau)\bmu=\frac{1}{L_\mu}\bmu,
\qquad
\mathbf{R}(\lambda,\tau)\bv=\frac{1}{B_v}\bv.
\]
Therefore
\[
g_{\mu\mu}
=
\bmu^\top \mathbf{R}(\lambda,\tau)\bmu
=
\frac{r}{L_\mu},
\]
\[
g_{\mu v}
=
\bmu^\top \mathbf{R}(\lambda,\tau)\bv
=
0,
\]
and, using the standing normalization \(\|\bv\|=1\),
\[
g_{vv}
=
\bv^\top \mathbf{R}(\lambda,\tau)\bv
=
\frac{1}{B_v}.
\]
Hence
\[
\Delta_G=g_{\mu\mu}g_{vv}-g_{\mu v}^2
=
\frac{r}{L_\mu B_v}.
\]

Substituting these identities into the denominator
\eqref{eq:lin_Dh} gives
\[
D_{\mathrm{proj}}(\alpha)
=
1+\tau\frac{r}{L_\mu}
+
\tau\phi\alpha^2\frac{1}{B_v}
+
\tau^2\phi(1-\phi)\alpha^2\frac{r}{L_\mu B_v}.
\]
Equivalently,
\[
D_{\mathrm{proj}}(\alpha)
=
\frac{
A_\mu B_v+\tau\phi P_\mu\alpha^2
}
{L_\mu B_v}.
\]
Thus, with
\[
D_{\mathrm{eig}}(\alpha):=A_\mu B_v+\tau\phi P_\mu\alpha^2,
\]
we have
\[
D_{\mathrm{proj}}(\alpha)
=
\frac{D_{\mathrm{eig}}(\alpha)}{L_\mu B_v}.
\]

Now use the exact projection formula \eqref{eq:lin_hmu_exact}. Since
\(g_{\mu v}=0\),
\[
h_\mu(\alpha)
=
\frac{
\tau\left((1-2\phi)\frac{r}{L_\mu}
+
\tau\phi(1-\phi)\alpha^2\frac{r}{L_\mu B_v}\right)
}
{
D_{\mathrm{proj}}(\alpha)
}.
\]
Multiplying numerator and denominator by \(L_\mu B_v\), we get
\[
h_\mu(\alpha)
=
r
\frac{
\tau(1-2\phi)B_v+\tau^2\phi(1-\phi)\alpha^2
}
{
D_{\mathrm{eig}}(\alpha)
}.
\]
This is the claimed formula for the benign projection.

Similarly, using \eqref{eq:lin_hv_exact} and \(g_{\mu v}=0\),
\[
h_v(\alpha)
=
\frac{
\tau\left(
\phi\alpha\frac{1}{B_v}
+
2\tau\phi(1-\phi)\alpha\frac{r}{L_\mu B_v}
\right)
}
{
D_{\mathrm{proj}}(\alpha)
}.
\]
Again multiplying numerator and denominator by \(L_\mu B_v\), we obtain
\[
h_v(\alpha)
=
\frac{
\tau\phi\alpha
\left(L_\mu+2\tau(1-\phi)r\right)
}
{
D_{\mathrm{eig}}(\alpha)
}
=
\frac{
\tau\phi Q_\mu\alpha
}
{
D_{\mathrm{eig}}(\alpha)
}.
\]
This proves the projection formulas in \eqref{eq:lin_ab_eigen_compact}.

It remains to prove the monotonicity and peak claims. Write
\[
D_{\mathrm{eig}}(\alpha)=D_0+D_2\alpha^2,
\qquad
D_0:=A_\mu B_v,
\qquad
D_2:=\tau\phi P_\mu .
\]
For the trigger projection,
\[
h_v(\alpha)
=
\frac{\tau\phi Q_\mu\alpha}{D_0+D_2\alpha^2}.
\]
Since \(\tau\phi Q_\mu>0\),
\[
h_v'(\alpha)
=
\frac{
\tau\phi Q_\mu(D_0-D_2\alpha^2)
}
{
(D_0+D_2\alpha^2)^2
}.
\]
Thus \(h_v\) is strictly increasing when \(\alpha^2<D_0/D_2\) and strictly
decreasing when \(\alpha^2>D_0/D_2\). Its unique maximizer satisfies
\[
(\alpha_\star^{\mathrm{eig}})^2
=
\frac{D_0}{D_2}
=
\frac{A_\mu B_v}{\tau\phi P_\mu},
\]
which is \eqref{eq:lin_alpha_star_eigen_compact}.

For the benign projection, write
\[
h_\mu(\alpha)
=
r\frac{N_0+N_2\alpha^2}{D_0+D_2\alpha^2},
\]
where
\[
N_0:=\tau(1-2\phi)B_v,
\qquad
N_2:=\tau^2\phi(1-\phi).
\]
Differentiating gives
\[
h_\mu'(\alpha)
=
r\frac{2\alpha(N_2D_0-N_0D_2)}
{(D_0+D_2\alpha^2)^2}.
\]
Using the definitions above,
\[
N_2D_0-N_0D_2
=
\tau^2\phi^2B_v
\bigl(L_\mu+2\tau(1-\phi)r\bigr)
=
\tau^2\phi^2B_vQ_\mu.
\]
This quantity is strictly positive. Therefore
\[
h_\mu'(\alpha)>0
\qquad
\text{for every } \alpha>0.
\]
This completes the proof of Corollary~\ref{cor:lin_peak_eigen}.
\end{proof}

\end{document}
